\subsection{Room/Mesa cohorts and paired experiment matrix}
\label{sec:room-design}

The separate room study uses the cleaned workbook, not the raw audit's
14\,649-row denominator. Its saved manifest records 14\,507 source rows,
14\,434 usable episodes and 73 exclusions, all zero clock endpoints in that
file version. Duration medians are learned from 10\,800 valid 2019--2021
cases, with a ten-observation threshold per procedure and a 70-minute global
fallback. Evaluation uses 2022. This temporal split permits recurring patients, unlike the patient-disjoint
model-preprocessing partitions described below.

The first room campaign contains one seven-case fixture, three historical
days (3 January: 9 cases; 2 February: 15; 20 October: 25), and two generated
samples at each of 100 and 300 cases. The selected dates have no exclusions; the middle and largest workloads
follow valid daily case counts among such dates. Generated instances sample complete eligible 2022
rows with replacement, retaining procedure--duration associations but not
within-day dependence. Each workload uses six methods and solver seeds
0, 1 and 2. Its 144 initial plans are each replayed under static/reactive
policies and closure/no-closure scenarios: 576 executions, not 576 independent
hospital days \citep{group2026roomstudy}.

Historical and generated days open 08:00--17:00 with 15-minute turnover.
Historical workloads use two interchangeable rooms; room 1 forbids starts
10:00--12:00 in the closure scenario. The small fixture uses two rooms,
08:00--11:00 and a 09:00--10:00 closure. Initial plans do not know the future
closure. For generated workloads,
\begin{equation}
 |\mathcal R|=\max\!\left(1,
 \left\lceil\frac{\sum_i(p_i+\tau)}{(b-a)\rho}\right\rceil\right),
 \qquad\rho=1.1.
\end{equation}
The four samples therefore use 14, 15, 44 and 45 rooms. Integer rounding
changes achieved load. This is an overloaded target experiment, not a
light-load sensitivity campaign; a single closure also represents a smaller
capacity fraction as the room pool grows.

Each nontrivial request has 1000 fitness calls and a 60-second safety limit.
No deadline stops occur in the saved campaigns. The baseline performs no
search evaluations. Reactive execution may use two additional request
budgets, so policies do not receive equal total computation. Runtime includes worker startup and communication. The 100-case runs were
sequential and 300-case runs used four concurrent batches, limiting timing
comparisons. Manifests record Python 3.12.3, Mesa 3.5.1, dependency versions
and source fingerprints, including dirty-checkout provenance.

\subsection{Duration-information experiment and predictive interfaces}
\label{sec:duration-design}

A separate saved campaign repeats the same matrix in median and oracle modes,
providing 1152 executions from 288 initial plans. These median runs are a
rerun, not extra independent replications of the first campaign. Oracle
mode replaces only scheduling estimates with realised room-occupancy minutes.
Workloads, execution durations and resources are held fixed; generated room
counts are computed from the original medians before substitution. Initial
plans are recomputed across modes but shared within each mode's four execution
conditions. Future closures remain hidden. This is a perfect-information
reference under finite search and a fixed decoder, not an ML predictor or a
certified operational upper bound. Normalised objective costs change scale
between modes; comparisons therefore use completions and realised overtime.

The operating-duration model branch \citep{group2026durationmodel} validates
patient-disjoint chronological partitions, excludes calendar features by
default, separates selection from later patient-disjoint calibration, and
selects the fastest candidate within 0.2 minutes of the best validation MAE.
Test data are reserved for reporting. Its review labels older metric tables
as belonging to a previous protocol; those are not verified current-model
results. Neither the median nor oracle campaign evaluates a learned room
model within Mesa.

Both preprocessing notebooks use pre-2021 training, new-patient validation
in 2021 and new-patient testing from 2022, excluding later cases of patients
in earlier partitions. Split identifiers/dates remain audit metadata.
Target-derived fields and postoperative/discharge information are excluded
from predictors; clinical codes require confirmed preoperative availability.
These controls must survive model-to-scheduler integration.

\subsection{Verification and reproducible evidence}

The incorporated room tables are regenerated from pinned aggregate CSVs.
Checks enforce unique run keys, four replays per initial plan, case conservation,
paired mode coverage and equality of no-closure policy outcomes. An independent
review of all 576 primary and 1152 duration execution exports recomputes
occupancy and overtime and checks start windows, closures, turnover separation,
evaluation caps and shared initial schedules. All pass, with zero reported
coordinator/search failures, timeouts or deadline stops. Validation of the reviewed source snapshot ran 94 tests, with 92 passing
and two optional tests skipped. These
checks support implementation consistency; they do not validate clinical
assumptions or make solver-seed repetitions independent observations.
